{
	\LARGE
	\flushleft \textit{\textbf{Author's Note}}
}
\addcontentsline{toc}{section}{Author's Note}\\
\vspace{1cm}

This entire note section is just about me and my musings, about this subject, and about me.
So, there isn't really much to learn here, this is simply my outlet regarding
what I have experienced and what I have to say.

I took Physics 131 about two semesters ago, when I was in fourth year. That was supposed
to be quite easy, Sir Paraan is quite lenient when it comes to subject outputs when compared
to other subjects, that being only giving a single problem set, and a singe exam. 
I aced the problem set and the exam during that time, having very high marks on both. However,
on the second requirement, that was where I was caught lacking. 

I never quite liked note taking during lectures. I believe it is distracting to a learner, and 
a student learns best while actively listening to the lecturer as he speaks, and digest these 
thoughts and learnings after the class is done, much like intellectual exercise.
As such, I never really took notes, not even in my highschool days, and simply learned everything 
through listening. This was how I worked 
for years, so taking notes was not in my nature as a student, and it worked for me, graduating 
high school distinguished, in fact.

The notebook was a pivotal
part of the subject, however, accounting for over 50\% of the grade, and so looking back
I really should have focused on doing it, much like a project. However, due to my rather
busy schedule that time, what with work and starting laboratory internships and other 
activities that took my time, I simply had no time to work on the notebooks. Looking back,
I should have just treated it the same way as weekly problem sets and should have done
my due diligence, but unfortunately, I was simply too swamped with workload that I couldn't
really work on it. I began the work only about five days before tha deadline, and couldn't 
complete it on time for said deadline. Futhermore, during the second submission, I was 
sufficiently burned out that I really couldn't do much anymore. When the time that I have
decided that I was not gonna pass this subject, the dropping deadline has lapsed, and so
I was stuck. It was at that time that I regretted not doing the notebook.

Anyhow, grading period came, and I recieved a 5.0, which is quite tragic, yes. It set me 
back another sem, and I should be graduating right now. However, it has passed, and it is 
no use to regret from our past mistakes, and instead take it as learning opportunities.

I am in fifth year now, a year more than one should have taken in their undergraduate
years. I'm academically old, about to graduate this year. I feel old, despite being 
twenty three. Despite my exemplary 
performance in my secondary education, I am significantly more laid back and 
exert the bare minimum effort during my college years. I am part of the pipeline 
of gifted students in their youth, turned academic disappointments during adulthood.
All those years seem wasted, but yes, I am now about to graduate from this god forsaken
University. It's time to get my shit together and get my act right.

So, for this semester, on account of having lesser workload in general, 
I have decided to go all out on the notebook, giving it my best effort, so that 
I may submit a great work when the times of submission comes, and I can be proud of it. 
Go big or go home, as they say.
Hopefully, what comes out of this work, though simply note taking and rather derivative, is of acceptable merit 
to both Sir Paraan, or his checker, and a work which, while I am averse in doing, I can 
look back on with mild happiness, and be proud.
\vfill
\begin{flushright}
    {\Large \textbf{End.}}
\end{flushright}
\begin{center}
	{\large
		\textsc{Physica Vincit Omnia}
	}
\end{center}
